\ssr{ЗАКЛЮЧЕНИЕ}

В результате работы были подтверждены различия в эффективности алгоритмов поиска в массиве. Бинарный поиск, обладающий сложностью $O(\log N)$, показал более высокую скорость по сравнению с методом полного перебора при работе с большими массивами. Это делает бинарный поиск предпочтительным выбором для работы с отсортированными массивами, где требуются высокие показатели производительности.

Метод полного перебора имеет сложность $O(N)$, что приводит к значительным затратам времени на выполнение поиска, особенно в худшем случае, когда элемент находится в конце массива или отсутствует. Таким образом, для неотсортированных массивов или в ситуациях, где простота реализации важнее эффективности, предпочтительнее использовать метод полного перебора, однако для оптимизации поиска в отсортированных данных бинарный поиск остается более эффективным вариантом.

В ходе выполненной работы были выполнены следующие задачи:
\begin{itemize}
    \item реализованы алгоритмы поиска значений в словаре;
    \item проведена оценка трудоемкости этих алгоритмов;
    \item выполнен сравнительный анализ сложности двух алгоритмов на основе замеров количества сравнений для каждого элемента массива и анализ данных для лучшего и худшего случаев.
\end{itemize}